%2025 - 10a
10a Rappel

$(M, \mc{O}, \mc{A})$

1. Dérivée covariante $\nabla$, $\nabla_XT$, $\to$ forme, $\to$ 2 tenseur

\[
\mc{T}^i_{(x)\ jk} = \Gamma^i_{(x)\ jk} \Gamma^i_{(x)\ kj} \qquad torsion\ (1,2)
\]
\[
\mc{R}^i_{jkl} = \partial\Gamma \cdot \Gamma .\Gamma \qquad courbure\ de\ Riemann\ (1,3)
\]

2. métrique g (0,2) symétrique inversible $g(\mc{V}_1,\mc{V}_2)=g(\mc{V}_2,\mc{V}_1)$

$T_PM$
\[
X\to Bra(X)
\]
"descente d'indice"
\[
Ket(\omega)\leftarrow \omega
\]
"montée d'indice"
\[
X=X^i\big(\frac{\partial}{\partial x^i}\big)
\]
\[
Bra(X)=\big[Bra(X)\big]_i\,dx^i
\]
\[
X_i\equiv\big(Bra(X)\big)_i=g_{ij}X^j\qquad g_{ij}
=g\big(\frac{\partial}{\partial x^i},\frac{\partial}{\partial x^j}\big)
\]
\[
\omega=\omega_idx^i\qquad Ket(\omega)=Ket(\omega)^i\frac{\partial}{\partial x^i}
\]
\[
\omega^i\equiv\big(Ket(\omega)\big)^i=h^{ij}\omega_j
\]
\[
\big[h^{ij}\big]=\big[g_{ij}\big]^{-1}\qquad h^{ij}g_{jk}=\delta^i_k
\]

géodésique "ligne droite"
\[
\boxed{
\begin{array}{ c }
 transport \\ parall\grave{e}le \\
\end{array}} \to ne\ tourne\ pas \quad \nabla_{\mc{V}_\gamma}\mc{V}_\gamma=0
\]

géodésique "chemin le plus court"
\[
L[\gamma]=\int^{\lambda_1}_{\lambda_2}
\underset{ \sqrt{g_{\mu\nu}(\dot x^\mu_\gamma,\dot x^\nu_\gamma)}}
{\sqrt{g(\mc{V}_\gamma,\mc{V}_\gamma)}}\,d\lambda
\]

\[
\delta L=0 \qquad \gamma\ \to\ x_\gamma\ \to\ x_\gamma + \delta x_\gamma
\]

\[
\nabla_{\mc{V}_\gamma}\mc{V}_\gamma=0 \quad\to\quad
\ddot x^i_\gamma+\Gamma^i_{jk}\ \dot x^j\ \dot x^k
\]
\[
L[\gamma]=\int\quad\to\quad\ddot x^i_\gamma+\overbrace{(g^{-1})^{ia}\big(
\frac{\partial}{\partial x^j}g_{k}
+\frac{\partial}{\partial x^k}g_{a}
-\frac{\partial}{\partial x^l}g_{l}\big)
}^{\displaystyle ^{L.C}\Gamma^i_{(x)\ jk}}
\ \dot x^j\ \dot x^k =0
\]

\chapter{Espace-temps physique et équations d'Einstein}
%Structure du cadre spatiotemporel : variété différentielle de dimension 4 munie d'une métrique Lorentzienne
%Contenu matériel : "points matériels", champs...
cadre, contenu, dynamique

\section{Structure}

\[
(\underset{\substack{\downarrow \\ dim\ 4}}{M}, \mc{O}, \mc{A}_\infty,
\underset{\substack{\downarrow \\ (1,3) \\ Lorentzienne}}{g}, ^{LC}\nabla_{(g)})
\]

cône de lumière

Minkowski g noté $\eta$
\[
\eta_{\mu\nu}=
\left[
  \begin{array}{cccc}
    -1 &  &  &  \\
     & 1 &  &  \\
     &  & 1 &  \\
     &  &  & 1 \\
  \end{array}
  \right]
\]
\[\]
\[\]
\[\]
\[\]
\[\]
